\section{概说和基本理念}
\begin{enumerate}
  \item 一个概念、两对关系、三组理念和六个范畴。
  \item 刑事诉讼是指在\zd{公检法专门机关}的主持下在\zd{当事人和其他诉讼参与人}的参加下，严格按照\zd{法定程序}，解决被追诉者的刑事责任问题的活动。
        \begin{enumerate}
          \item \zd{3 大 4 小来主持}。
          \item \zd{诉讼参与人参加} = 当事人 + 其他诉讼参与人。
          \item 程序。立案 \yjt 侦查 \yjt 起诉 \yjt 审判 \yjt 执行。
        \end{enumerate}
  \item 刑诉法。调整刑诉活动的法律规范的总称。有狭义（刑法典）、广义（调整刑诉活动的一切法律规范）。
  \item 渊源。在哪些文件中能找到刑诉法的依据。判断标准。\zd{以下 7 个 + 内容相关}。
        \begin{enumerate}
          \item 宪法、刑诉法典。
          \item 法律和立法解释、法律解释；规定、法规；部门规章。
          \item 国际条约和公约。
        \end{enumerate}
  \item 刑诉法和刑法的关系。
        \begin{enumerate}
          \item \zd{工具价值}。\zd{提供、保障}。为刑法提供组织、基本构架、程序证据规范、标准、准确性、效率性上的保障。
          \item \zd{独立价值}。\zd{本身、弥补、影响}。
                \begin{enumerate}
                  \item 本身体现民主法制。
                  \item 弥补刑法不足/空白。实体对实体。\al{特别追诉}。
                  \item 影响刑法结果实现。程序对实体。\al{刑事和解、排非、上诉不加刑}。\al{特别追诉 \cuo}。
                \end{enumerate}
        \end{enumerate}
  \item 与法治国家的关系。宪法是\zd{静态}的刑诉法，刑诉法是\zd{动态}的宪法。宪法不会自我实现它的内容，刑诉法进行维护捍卫。
  \item 基本理念。刑事诉讼的灵魂。整个刑诉流程以及各项制度设计的指导思想。
        \begin{enumerate}
          \item 惩罚犯罪。\zd{打击犯罪}。
          \item 保障人权。
                \begin{enumerate}
                  \item \zd{无辜的人}不受追究、\zd{有罪的人}受到公正处罚、\zd{诉讼权利}得到充分保障和行使。
                  \item 制度涉及。辩护权、法院援助、非法证据排除。
                \end{enumerate}
          \item \zd{保障人权和惩罚犯罪既对立又统一，二者要并重}。
          \item 实体公正。\zd{结果}公正。事实清楚、证据充分；定罪准确、量刑适当；涉案财物处置公正、有错必纠。
          \item 程序公正。\zd{过程}公正。程序要遵守；权利要保障；措施要适当；刑讯要严禁；程序要透明；审判要中立、控审要分离；期限要及时。
          \item \zd{坚持实体公正和程序公正并重 4 原则}。
          \item 司法公正。公正是\zd{首要目标价值}。
          \item 诉讼效率。速度、快捷、节省时间、办案效率。简易程序、速裁程序、缺席审判、认罪认罚。
                \begin{enumerate}
                  \item \al{对未成年附条件不起诉 \dui{起诉便宜主义}}。
                  \item \al{交通肇事，不逮捕而直接起诉 \dui{办案效率}}。
                  \item \al{为保证证人安全，庭前核实身份 \dui{节省庭审时间}}。
                  \item \al{重新开庭就量刑部分组织讨论 \cuo{为了公正牺牲效率}}。
                \end{enumerate}
          \item \zd{公正优先，兼顾效率}。
        \end{enumerate}
  \item 量词。以上、以下、以内、届满包括本数。不满、超过、以外，不包括本数。
\end{enumerate}
